\section{Coloring}

\begin{definition}

    A \textbf{$k$-coloring} of $G = (V, E)$ is a mapping $c: V \to [k]$ such that $c(u) \neq c(v)$ for all
    $\{u,v\} \in E$.
    \\

    The \textbf{chromatic number} $\chi(G)$ is the minimum number of colors.
    \\

    $\chi(G) \le 2 \iff G$ is bipartite.
    
\end{definition}

\begin{theorem}

    For $k \ge 3$, the problem "Is $\chi(G) \le k$?" is NP-complete.
    
\end{theorem}

\begin{algorithm}

#Greedy Coloring
#Requires at most maximum degree + 1 colors.

Heuristic: 
delete vertices in order of minimum degree
#Gives <= k+1 colors if all subgraphs have a vertex with degree <= k.
    
\end{algorithm}

\begin{theorem}

\textbf{Brooks' Theorem}:\\

    For a connected graph $G$ that is neither a complete graph $K_n$ nor an odd cycle $C_{2n+1}$:
    $$\chi(G) \le \Delta(G)$$
    
\end{theorem}

